\scanpage{40}
\margin{\(^*\) Il suffit \unsure{pour que l'on ait} que
\(\delta=d\cap d'\) soit segment initial de \(d'\) et contienne
le plus grand élément de \(d\), ce qui dit que pour les
\(\textcircled{D}\), \(\Sigma_d-\Sigma_\delta\) est
\unsure{mono.} et \(\Sigma_{d'}-\Sigma_d\) est fibrant.
\unsure{Dans plus}}

(3) Soient \(d,d'\) deux drapeaux tels que
le dernier élément de \(d\) soit le premier élément de \(d'\)
et considérons \(d''=d\cup d'\) et le diagramme
\diagram{\begin{tikzcd}
 \Sigma_{d''}\arrow[r,hook]\arrow[d]&\Sigma_{d'}\arrow[d]\\
 \Sigma_d\arrow[r,hook]&\Sigma_i
\end{tikzcd}}
\unsure{est}
est \uline{cartésien}. (Pour \(d\) ou \(d'\)
de longueur \(0\) i.e. \unsure{n'étant} : \(\{i\}\), cette condition
est triviale.)

\margin{Il suffit de le supposer pour
\(d\) de longueur \(2\) (ou \(d'\) de longueur \(2\)).
On en déduit plus
\deleted{\unsure{la généralité que si}}
\deleted{\unsure{quelconques}}
\(d,d'\) sont deux drapeaux tels que
\inserted{\(d''=d\cup d'\)}
soit un drapeau et que
\(\delta=d\cap d'\)
\deleted{\(d\cup d'\)}
\(d''=d\amalg_\delta d'\)
\unsure{cat. ou un} \deleted{\unsure{ens}} ordonné,
\(^*\) alors
\[
 \textcircled{D}
\]
\diagram{\begin{tikzcd}
 \Sigma_{d''}\arrow[r]\arrow[d]&\Sigma_{d'}\arrow[d]\\
 \Sigma_d\arrow[r]&\Sigma_\delta
\end{tikzcd}}
est cartésien.}

\margin{De plus, il faut supposer que lorsque
\(d\) est de longueur \(2\), alors
\([\Sigma_{d'}\unsure{-|\longrightarrow}\Sigma_\delta]\) est
\uline{transverse} aux sous-espaces cylindriques \(\Sigma_d\) de
\(\Sigma_{\unsure{\delta}}\) --- et idem pour les structures cylindriques.}

\margin{(\unsure{comme au n° précédent})
On suppose de
plus que \deleted{\unsure{d'ailleurs}} si pour \(i\in I\),
\(\operatorname{Dr}(I,i)\) désigne le sous-ens ordonné
de \(\operatorname{Dr}(I)\) formé des drapeaux
\(d\) dont le dernier élément est \(i\),
de sorte que les \(\Sigma_d\) sont \(\subset\Sigma_i\),
alors les \(\Sigma_d\) forment dans \(\Sigma_i\) une
\uline{stratification précylindrique} de
type \(\operatorname{Dr}(I,i)^\circ\), ce qui s'exprime par
les conditions supplémentaires
\[
 \text{a)}\quad \Sigma_{kj}\cap\Sigma_{\ell j}=\varnothing
 \quad\text{si }k,\ell\unsure{\le}j\text{ sont des noncomparables}.
\]
\[
 \text{b)}\quad
 \Sigma_{i_0\ldots i_n}
 =\Sigma_{i_0i_n}\cap\Sigma_{i_1i_n}\cdots
 \cap\Sigma_{i_{n-1}i_n}
 \qquad\text{dans }\Sigma_{i_n}
\]
\[
 \left(\Longleftrightarrow\
 \Sigma_d\cap\Sigma_{d'}=
 \begin{cases}
 \varnothing&\text{si }d\cup d'\text{ \unsure{pas un drapeau}}\\
 \Sigma_{d\cup d'}&\text{sinon}
 \end{cases}\right)
\]
c) L'intersection dans b) est transversale.
On suppose plus précisément, \unsure{choisis|choisies} des structures
cylindriques \unsure{qui rendent transversales}
\(\Sigma_{ij}\) \(\unsure{(i<j)}\).}

Cela nous donne un procédé \unsure{canonique} pour
reconstruire les \(\Sigma_d\), à isom. \inserted{can.} près,
quand on connaît les \(\Sigma_i,\Sigma_{ij}\)
\inserted{\((i<j)\)} et les morphismes
\diagram{\begin{tikzcd}
 \Sigma_{ij}\arrow[r,hook]\arrow[d]&\Sigma_j\\
 \Sigma_i&
\end{tikzcd}}
\deleted{Pour}
(NB \unsure{que}
la flèche verticale est fibrante, et l'horizontale
un morphisme-bord). Pour avoir les
morphismes de transition entre les
\(\Sigma_d\) (construits par récurrence sur \(\operatorname{long}d\),
les \(0\) et \(1\) étant déjà donné) il faut se donner
les morphismes
\[
 \underbrace{\Sigma_{ij}\mathbin{\times_{\Sigma_j}}\Sigma_{jk}}
 _{\mathrel{\mathop=^{\mathrm{déf}}}\Sigma_{ijk}}
 \longrightarrow\Sigma_{i,k}
\]
satisfaisant aux conditions

\margin{39}
